\section{Results}
\label{sec:results}

We present a comprehensive evaluation of our multimodal framework across two prediction tasks: gene expression and cell type composition prediction. Multimodal fusion consistently outperforms unimodal baselines, spatially aligned token-level transcript representations are crucial for capturing tissue heterogeneity, and lightweight task-specific encoders can outperform large pretrained foundation models.

\paragraph{Multimodal Fusion Outperforms Unimodal Baselines}
\label{par:multimodal-fusion-outperforms-unimodal-baselines}

Multimodal fusion models outperform all unimodal baselines on both prediction tasks (\autoref{fig:p-fus:gene-lvl_2}, \autoref{fig:p-fus:cell-lvl_2}), with full results reported in \autoref{tab:p-fus:gene-table} and \autoref{tab:p-fus:cell-table}. For gene expression prediction, our best multimodal models achieve a $26.5\%$ improvement in Spearman correlation over the best morphology-only baseline and a $13\%$ improvement over the best expression-only baseline, demonstrating the complementary nature of the two modalities.

\paragraph{Superior Performance on Cell Type Count Prediction}
\label{par:superior-performance-on-cell-type-count-prediction}

Multimodal fusion achieves even stronger gains on cell type composition prediction, substantially outperforming all other configurations (\autoref{tab:p-fus:cell-table}, \autoref{fig:p-fus:cell-lvl_2}). Cell type composition prediction is where the early-add fusion approach most clearly differentiates itself, consistently outperforming all other model variants (\autoref{tab:p-fus:cell-table}, \autoref{fig:p-fus:cell-lvl_2}). Notably, unlike the gene expression task where late fusion remains competitive, early-add fusion dominates the late fusion approach here. We attribute this to the inherent difficulty of the task: accurately reproducing cellular heterogeneity and spatial organization requires disentangling distinct phenotypes and their relative abundances, a task that benefits substantially from the fine-grained integration of morphological and molecular features afforded by early fusion.

\paragraph{Token-Level Transcript Representations Capture Spatial Heterogeneity}
\label{par:token-level-transcript-representations-capture-spa}

Token-level transcript representations are a key driver of performance for gene expression prediction. Token-level early fusion slightly outperforms \emph{vanilla} late fusion with tile-level transcripts. However, the decisive factor is not fusion timing \textit{per se} but the granularity of the transcript representation: when late fusion is augmented with spatially aligned, token-level transcript pooling, the performance gap closes, and tokenized late fusion performs on par with early fusion.

For cell type prediction, tokenized late fusion continues to outperform \emph{vanilla} late fusion even in the absence of this advantage in expression-only models (\autoref{tab:p-fus:cell-table}, \autoref{fig:p-fus:cell-lvl_2}), underscoring the benefits of token-level representations.

To further validate the value of token-level aggregation, CLIP-style contrastive pretraining (\autoref{tab:p-fus:clip-results}) confirms that token-level aggregation improves cross-modal retrieval over tile-level aggregation (recall@1: $0.062 \pm 0.004$ vs.\ $0.041 \pm 0.003$; random baseline: $0.004$), demonstrating that spatially resolved transcript features yield better-aligned multimodal embeddings.

\paragraph{Task-Specific Training Outperforms Frozen Foundation Models}
\label{par:task-specific-training-outperforms-frozen-foundati}
\leavevmode\\ Lightweight, task-specific encoders consistently outperform frozen H\&E foundation model embeddings. A low-capacity ViT-S trained end-to-end with a simple MLP expression encoder achieves higher correlations than linear probes on frozen CONCH, Phikon, or OmiCLIP/Loki features (\autoref{tab:p-fus:gene-table}), suggesting that task-specific adaptation to the morpho-molecular relationships relevant to spatial transcriptomics prediction can outweigh the benefits of large-scale pretraining.

\paragraph{Performance Analysis and Spatial Predictions}
\label{par:performance-analysis-and-spatial-predictions}

Both early and late fusion models clearly outperform the Min-MAE baseline across all genes and functional modules (\autoref{fig:p-fus:main-results-minmae}). The expression-only model, by contrast, performs comparably to Min-MAE, indicating that its predictions are largely driven by simple cross-panel gene--gene correlations (\autoref{fig:p-fus:src-tgt-corr}, \autoref{fig:p-fus:avg-src-corr}). The morphology-only model is outperformed by Min-MAE across almost all genes, confirming that morphology alone provides limited information for predicting specific gene expression values without transcriptomic context.

Spatially, our predictions closely mirror the true measurements (\autoref{fig:p-fus:true-vs-predicted}, \autoref{fig:p-fus:cell-type-spatial-pred}), capturing rich morpho-molecular structures of the tumor microenvironment (\autoref{fig:p-fus:umaps}), indicating that the models learn biologically meaningful representations of tissue architecture and cellular states.

\begin{figure}[t]
  \centering
  \begin{subfigure}[b]{\textwidth}
    \centering
    \includegraphics[width=0.65\textwidth]{cell-boxplot-marsilea.pdf}
    \caption{\textbf{Cell type composition prediction.}
      Pearson correlation coefficient across all 39 cell types for each tested model configuration.
      Multimodal fusion consistently outperforms unimodal baselines, with early fusion achieving the strongest gains.}
    \label{fig:p-fus:cell-lvl_2}
  \end{subfigure}

  \vspace{1em}

  \begin{subfigure}[b]{\textwidth}
    \centering
    \includegraphics[width=0.55\textwidth]{gene-pcc-boxplot-marsilea.pdf}
    \caption{\textbf{Gene expression prediction.}
      Spearman correlation coefficient across all add-on panel genes for each tested model configuration.
      Multimodal fusion consistently outperforms unimodal baselines, demonstrating the complementary nature of morphological and transcriptomic features.}
    \label{fig:p-fus:gene-lvl_2}
  \end{subfigure}

  \caption{\textbf{Prediction performance across all model configurations.}
    Multimodal fusion consistently outperforms unimodal baselines on both tasks.}
  \label{fig:p-fus:combined}
\end{figure}


\begin{figure}[t]
  \centering
  \includegraphics[width=0.95\textwidth]{cell-all-pcc-boxplot-marsilea.pdf}
  \caption{\textbf{Cell type prediction performance stratified by input gene panel.}
  Pearson correlation for models trained on the full panel, the core panel, and the add-on panel as transcript input.
  Core-panel models closely match full-panel performance, while add-on-only models degrade substantially,
  with multimodal integration providing the largest relative benefit in this low-information regime.}
  \label{fig:p-fus:cell-panel}
\end{figure}

\begin{table}
\centering
\caption{\textbf{CLIP pretraining retrieval performance.} Image--transcript retrieval results (recall@1 and recall@5; mean$\pm$std across runs) for different core-transcript aggregation strategies. Random baselines correspond to chance performance at batch size 256.}
\label{tab:p-fus:clip-results}
\begin{tabular}{lllll}
\toprule
expr\_encoder & morph\_encoder & pool & recall@1 & recall@5 \\
\midrule
mlp & vit\_s & token & 0.062±0.004 & 0.202±0.009 \\
mlp & vit\_s & tile  & 0.041±0.003 & 0.151±0.011 \\
\midrule
random & -- & -- & 0.004 & 0.020 \\
\bottomrule
\end{tabular}
\end{table}

\paragraph{Model Capacity Matters}
\label{par:model-capacity}

We compare low-capacity configurations (ViT-S with a single-layer, 32-unit MLP expression encoder) against high-capacity configurations (ViT-B with a two-layer, 128-unit MLP expression encoder) across all fusion strategies and stages for the cell type prediction task (\autoref{tab:p-fus:cell-table}, \autoref{fig:p-fus:high-low-capacity}).

Increasing model capacity consistently improves multimodal and expression-only models but yields no gains for morphology-only baselines. For early-add fusion, the high-capacity ViT-B achieves a mean Pearson correlation of $0.761$ versus $0.688$ for ViT-S, while morphology-only ViT-B and ViT-S perform comparably ($0.554$ vs.\ $0.541$--$0.543$). For expression-only models, increasing MLP depth and width improves performance from $0.576$ to $0.655$. These results motivate larger architectures in the multimodal setting, where richer inputs can be more fully exploited.

\paragraph{Core Panel Input Is Sufficient}
\label{sec:panel-cell-type}

We evaluate models trained on either \emph{all} available genes, the \emph{core} panel alone, or the \emph{add-on} panel alone as transcript input for cell type prediction (\autoref{fig:p-fus:cell-panel}). Models with access to the full or core panel substantially outperform those using only the add-on panel, while the gap between full- and core-panel models is modest. This is consistent with the biological design of the panels: the core Human Immuno-Oncology Panel covers established immune and cell-type-specific markers directly informative for cell identity, whereas the add-on panel contains study-specific probes.

Multimodal models consistently outperform their expression-only counterparts regardless of input panel, confirming that morphological context provides complementary, non-redundant information. The morphological advantage is largest when only the add-on panel is available.

\paragraph{Per-Cell-Type Analysis Reveals Characteristic Unimodal Failure Modes}
\label{sec:per-cell-type}

Aggregate correlation statistics can mask important per-cell-type variation. \autoref{fig:p-fus:sample-level-pearson} summarizes the distribution of Pearson correlations across test splits for each of the 39 cell types (ordered left to right by decreasing abundance) as a boxplot, while \autoref{fig:p-fus:sample-level-pearson-2} shows sample-level correlations for a randomly selected subset of samples as a dotplot. Both figures show that performance is strong and consistent for the majority of abundant cell types, with a modest drop toward rarer populations on the right. To further investigate cell type specific performances across models, we highlight encoder-specific collapse modes, for which we focus on a representative filtered subset (\autoref{fig:p-fus:cell-type-level}) and compare early-add, expression-only, and morphology-only directly.

Early-add fusion achieves high and consistent correlations across the majority of cell types. Unimodal models, by contrast, exhibit characteristic failure modes that fall into two distinct patterns.

The first pattern occurs when both unimodal models perform poorly, yet their combination substantially outperforms either alone. Regulatory T cells are the most extreme case: expression-only correlations are low, morphology-only is only marginally better, yet early-add achieves substantially higher correlations. The same pattern holds for ciliated epithelial, club, and mesothelial cells.

The second pattern occurs when expression-only specifically fails while morphology-only retains reasonable signal, as seen in natural killer cells and type 1 interferon response T cells. The effect is most striking for plasmacytoid dendritic cells, where both early-add and morphology-only reach near-perfect correlation (${\sim}0.95$) while expression-only collapses to around $0.60$, indicating that their identity is strongly encoded in morphology but poorly represented in the core transcriptomic panel.

Together, these patterns provide a mechanistic rationale for early-add fusion that goes beyond aggregate performance, with the nature of the gain varying across cell types.

\begin{figure}
    \centering
  \includegraphics[width=\textwidth]{cell_type_level_pearson_boxplot-early-only-all-cells.pdf}
\caption{\textbf{Per-cell-type Pearson correlation for the early-add model across all 39 cell types.} Cell types are ordered by decreasing abundance from left to right. Correlations are high and consistent for the majority of abundant cell types, with 14 of the top 15 most abundant populations exceeding 0.8. Performance declines for rarer populations toward the right, with regulatory T showing the lowest, collapse-like correlations.}
    \label{fig:p-fus:sample-level-pearson}
\end{figure}

\begin{figure}[htbp]
  \centering
  \includegraphics[width=\textwidth]{cell_type_level_pearson_boxplot-target-models-filtered.pdf}
\caption{\textbf{Per-cell-type Pearson correlation for early-add and unimodal configurations on a representative filtered subset of cell types.} Boxes show the distribution of Pearson correlations across test splits for early-add (orange), expression-only (blue), and morphology-only (green). Early-add correlations are consistently higher across all cell types. Morph-only correlations collapse markedly for plasmacytoid dendritic cells, expr-only for natural killer cells. Unimodal correlations are markedly lower for regulatory T, ciliated epithelial, club, and mesothelial cells, compared to the fusion variant.}
  \label{fig:p-fus:cell-type-level}
\end{figure}

\begin{figure}
    \centering
    \includegraphics[width=0.5\linewidth]{XE_1G11.03_HNE_1G1103.pdf}
    \caption{\textbf{Spatial cell type abundance predictions for a representative sample.}
    Predicted versus true cell type abundances mapped onto the tissue section
    for selected cell types, demonstrating that the multimodal model captures
    spatially coherent patterns of cellular composition.}
    \label{fig:p-fus:cell-type-spatial-pred}
\end{figure}

% \begin{figure}
%     \centering
%     \includegraphics[width=0.5\linewidth]{papers/Fusion/figures/cell-viz/bs6thp97/XE_1FXS.01_HNE_1FXS01.pdf.png}
%     \caption{Cell type abundance prediction for selected cell types.}
%     \label{fig:placeholder}
% \end{figure}

\begin{figure}[t]
  \centering
  \includegraphics[width=\linewidth]{UMAPs.pdf}
  \caption{\textbf{Learned embeddings capture rich morpho-molecular descriptions of the tumor microenvironment.} (\textbf{A}) Leiden clustering of image/tile embeddings for a single sample across models. For fairness, embeddings are reduced by PCA to 31 dimensions (matching the MLP bottleneck capacity; MLP uses 32) and clustered with fixed Leiden resolutions (0.5 and 1.0) without per-model tuning. Expr-only yields granular, noisy clusters reflecting fine-grained transcriptomic variation; Morph-only produces larger coherent regions with potential under-clustering; early fusion better reflects transcript-driven heterogeneity while remaining spatially coherent; late fusion closely mirrors Morph-only patterns. Cluster colors are not aligned across methods.
(\textbf{B}) Corresponding UMAP visualizations of the embeddings colored by the Leiden clusters from (A), summarizing global neighborhood structure under identical dimensionality reduction and clustering settings. The UMAP projections reveal how different fusion strategies organize tissue regions in embedding space.}
  \label{fig:p-fus:umaps}
\end{figure}

\paragraph{BEAT Evaluation with Highly Variable Gene Targets}

\paragraph{Sample-Level Performance}

\paragraph{Sample-Level Variance}

\paragraph{Gene- and Cell-Type-Level Performance}

\paragraph{Gene- and Cell-Type-Level Variance}

\paragraph{Representation Geometry of Learned Embeddings}

\paragraph{MIL Attention Patterns}
